La interfaz permite a la persona usuaria iniciar el entrenamiento de nuevos modelos utilizando el mismo pipeline empleado durante el desarrollo experimental del proyecto. Los parámetros definidos en la sección de configuración avanzada se utilizan para construir la configuración del entrenamiento y lanzar un nuevo ciclo de optimización. Los checkpoints generados por estos entrenamientos se almacenan de forma independiente a los modelos utilizados en los experimentos principales.

Esta funcionalidad tiene principalmente un objetivo demostrativo. Debido al coste computacional del entrenamiento, no se pretende que una persona usuaria obtenga mediante la interfaz un modelo competitivo de manera inmediata, sino permitirle observar el proceso de entrenamiento y experimentar con determinados hiperparámetros para comprender cómo afectan al funcionamiento del sistema.
Como el proceso de entrenamiento requiere de un gran tiempo de espera obtener un modelo entrenado, seguramente la persona usuaria no llegue a esperar lo suficiente para haber realizado un entrenamiento que ha aprendido lo suficiente como para generar resultados interesantes, se ha añadido esta funcionalidad porque se 
considera interesante que el persona usuaria pueda ver en tiempo real cómo es realmente un entrenamiento de un modelo de inteligencia artificial.